We start our discussion from Lemma~\eqref{lem:forward-newtons-formula-for-powers},
which yields forward Newton's interpolation formula for powers.
For example,
\begin{align*}
  n^3 &= 0 \tbinom{n-0}{0}  + 1 \tbinom{n-0}{1}  + 6 \tbinom{n-0}{2}  + 6 \tbinom{n-0}{3}, \\
  n^3 &= 1 \tbinom{n-1}{0}  + 7 \tbinom{n-1}{1}  + 12 \tbinom{n-1}{2} + 6 \tbinom{n-1}{3}, \\
  n^3 &= 8 \tbinom{n-2}{0}  + 19 \tbinom{n-2}{1} + 18 \tbinom{n-2}{2} + 6 \tbinom{n-2}{3}, \\
  n^3 &= 27 \tbinom{n-3}{0} + 37 \tbinom{n-3}{1} + 24 \tbinom{n-3}{2} + 6 \tbinom{n-3}{3}.
\end{align*}
Thus, the transition to formula for sums of powers is quite straightforward,
% Hence, by Newton's formula for
% powers~\eqref{prop:newtons-interpolation-formula-for-powers},
% we get formula for ordinary sums of powers.
\begin{align*}
  \KnuthRFoldSum{1}{n}{m}
  = \sum_{k=0}^{m} \left[ \sum_{j=1}^{n} \binom{j-t}{k} \right] \Delta^{k} t^m.
\end{align*}
Therefore, by setting $r=0$ into Lemma~\eqref{lem:forward-hockey-stick-identity}
yields closed form of ordinary sums of powers
\begin{proposition}[Ordinary sums of powers]
  \label{prop:ordinary-sums-of-powers}
  For non-negative integers $n,m$, and an arbitrary integer $t$
  \begin{align*}
    \KnuthRFoldSum{1}{n}{m}
    = \sum_{k=0}^{m} \left[ \binom{n-t+1}{k+1} - \binom{1-t}{k+1} \right] \Delta^{k} t^m.
  \end{align*}
\end{proposition}
We may see that the basis $\binom{j-t+r}{k+r}$ for repeatedly
applied Lemma~\eqref{lem:forward-hockey-stick-identity}
was indeed chosen correctly, because we derived closed form for sums of powers
with ease, simply setting parameter for $r$.
For example,
\input{examples/01_ordinary_sums_of_cubes.tex}
From example above, we may observe binomial coefficients of negative upper
index, which may looking unfamiliar.
However, there is no reason to be afraid of, because negative
upper index binomial coefficients are completely deterministic,
because $\binom{-n}{k} = (-1)^{k} \binom{k+n-1}{k}$.
By setting $t=0$ into
Proposition~\eqref{prop:ordinary-sums-of-powers}
yields one more well-known identity for sums of powers,
mentioned in~\cite[p. 190]{graham1994concrete}, and~\cite{pfaff2007deriving}
\begin{corollary}
  For non-negative integers $n,m$
  \begin{align*}
    \KnuthRFoldSum{1}{n}{m} = \sum_{j=0}^{m} \binom{n+1}{j+1} \Delta^{j} 0^m.
  \end{align*}
\end{corollary}
For example,
\input{examples/02_ordinary_sums_of_powers_t_0.tex}
From the example above, we observe the coefficients $0,1,2,0, 1, 6, 6, 0, 1, 14, 36, 24, \dots$,
which is the sequence
\href{https://oeis.org/A131689}{\texttt{A131689}}
in the OEIS~\cite{sloane2003line}.
% where $t$ is an arbitrary integer.
% By shifted hockey stick
% identity~\eqref{lem:shifted-generalized-hockey-stick-identity},
% the closed form of ordinary sums of powers yields.

% \begin{proposition}[Ordinary power sums]
%   \label{prop:ordinary-sums-of-powers-via-newtons-formula}
%   For integers $n\geq 0$, and $m\geq 0$, and an arbitrary integer $t$,
%   \begin{align*}
%     \KnuthRFoldSum{1}{n}{m} = \sum_{j=0}^{m} \Delta^{j} t^m
%     \left[ \binom{n-t+1}{j+1} - \binom{1-t}{j+1} \right].
%   \end{align*}
% \end{proposition}
